\documentclass{article}

\usepackage{fullpage}
\usepackage{url}
\usepackage{graphicx}
\usepackage{amsmath}
\usepackage{physics}
\usepackage{mathtools}
\usepackage{bbold}

\DeclareMathOperator{\erfi}{erfi}
\newcommand{\R}{\ensuremath{\mathbb{R}}}

\title{Quantum Mechanics II Excercise 3}
\author{Idan Stark}
\date{\today}

\begin{document}
\maketitle

\section{$SU(2), SO(3)$}
\subsection{$SU(2)$ and the Pauli matrices}
$SU(2)$ is defined as the set of all $2\times2$ complex matrices such that $U^\dagger U = \mathbb{1}$ and $\det U = 1$. We can characterise those matrices by starting from a generic $2\times2$ matrix and narrowing down using the conditions.
\newline
Given a $2\times2$ matrix U, it can be written as
\begin{equation}
    U = \begin{pmatrix}a & b \\ c & d\end{pmatrix}
\end{equation}
The first condition for it to be in $SU(2)$ is $U^\dagger U = \mathbb{1}$. This means that the matrix is invertible and that its inverse is its complex conjugate
\begin{equation}
    U^\dagger = U^{-1}
\end{equation}
Using the formula for the inverse of a $2\times2$ matrix:
\begin{equation}
    U^{-1} = \frac{1}{|\det U|}\begin{pmatrix} d & -b \\ -c & a \end{pmatrix}
\end{equation}
And using the second condition $\det U = 1$ we get:
\begin{equation}
    \begin{pmatrix} a^* & c^* \\ b^* & d^* \end{pmatrix}
    =
    \begin{pmatrix} d & -b \\ -c & a \end{pmatrix}
\end{equation}
\begin{equation}
    a^* = d \qquad d^* = a \qquad b^* = -c \qquad c^* = -b
\end{equation}
In other words, we can write $U$ as:
\begin{equation}
\label{eq:su_complex_numebrs_rep}
    U = \begin{pmatrix} a & b \\ -b^* & a^* \end{pmatrix}
\end{equation}
Writing $a$ and $b$ in their complex cartesian representation, we get:
\begin{equation}
    \label{eq:final_u}
    U = \begin{pmatrix} x + iy & w + iz \\ -w + iz & x - iy \end{pmatrix}
\end{equation}
And requiring $\det U = 1$ gives us:
\begin{equation}
    \label{eq:final_u_condition}
    x^2 + y^2 + w^2 + z^2 = 1
\end{equation}
Now, let us look at the expression
\begin{equation}
    A = \exp\left[i t_i \sigma_i\right]
\end{equation}
Using the definition of the Pauli matrices, we get
\begin{equation}
    A = \exp\left[\frac{1}{2}i\theta_i \sigma_i\right] = \exp
        \left[\frac{1}{2}\begin{pmatrix}
            i\theta_z & \theta_y + i\theta_x \\
            - \theta_y+ i\theta_x & -i\theta_z
        \end{pmatrix}\right] \equiv \exp \left[\frac{1}{2}P(\theta)\right]
\end{equation}
Where we named the matrix in the exponent $P(\theta)$. We can now expand the exponent, but first let us calculate the powers of $P(\theta)$:
\begin{equation}
    P^2(\theta) = \begin{pmatrix}
            i\theta_z & \theta_y + i\theta_x \\
            - \theta_y+ i\theta_x & -i\theta_z
        \end{pmatrix}\begin{pmatrix}
            i\theta_z & \theta_y + i\theta_x \\
            - \theta_y+ i\theta_x & -i\theta_z
        \end{pmatrix}
        =
        -\begin{pmatrix}
            \theta_x^2 + \theta_y^2 + \theta_z^2 & 0 \\ 0 & \theta_x^2 + \theta_y^2 + \theta_z^2
        \end{pmatrix} \equiv -\theta^2 \mathbb{1}
\end{equation}
Therefore, we can split the powers of $P$ to even and odd powers:
\begin{equation}
    P^{2n}(\theta) = \left(-1\right)^n\theta^{2n}\mathbb{1} \qquad P^{2n+1}(\theta) = \left(-1\right)^n\theta^{2n}P
\end{equation}
And thus
\begin{equation}
\begin{gathered}
\label{eq:su_final_rep}
    A = \exp \left[\frac{1}{2}P\right] = \sum_{n=0}^\infty \frac{(P/2)^n}{n!} =
    \mathbb{1}\sum_{n \ even}\frac{\left(i\theta/2\right)^n}{n!} + \frac{1}{2}P\sum_{n \ odd}\frac{\left(i\theta/2\right)^{n - 1}}{n!} = \cos(\frac{\theta}{2})\mathbb{1} + \sin(\frac{\theta}{2}) \frac{P}{\theta} = \\
    = \begin{pmatrix}
        \cos(\theta/2) + i\sin(\theta/2)\frac{\theta_z}{\theta} & \sin(\theta/2)\frac{\theta_y + i\theta_x}{\theta} \\ 
        \sin(\theta/2)\frac{-\theta_y + i\theta_x}{\theta} & \cos(\theta/2) - i\sin(\theta/2)\frac{\theta_z}{\theta}
    \end{pmatrix}
\end{gathered}
\end{equation}
The determinant of this matrix is:
\begin{equation}
\begin{gathered}
    \det A = \left(\cos(\theta/2) + i\sin(\theta/2)\frac{\theta_z}{\theta}\right) \left(\cos(\theta/2) - i\sin(\theta/2)\frac{\theta_z}{\theta}\right) - \sin^2(\theta/2)\frac{- \theta_x^2 - \theta_y^2}{\theta^2} = 
    \\
    = \cos^2(\theta/2) + \sin^2(\theta/2)\frac{\theta_z^2}{\theta^2} + \sin^2(\theta/2)\frac{\theta_x^2 + \theta_y^2}{\theta^2} = \cos^2(\theta/2) + \sin^2(\theta/2)\frac{\theta_x^2 + \theta_y^2 + \theta_z^2}{\theta^2} =
    \\
    = \cos^2(\theta/2) + \sin^2(\theta/2) = 1
\end{gathered} 
\end{equation}
Which means that with a small change of variables we can make it fit the matrix in equation \ref{eq:final_u}:
\begin{equation}
    x = \cos(\theta/2) \qquad y = \sin(\theta/2)\frac{\theta_z}{\theta}
    \qquad w = \sin(\theta/2)\frac{\theta_y}{\theta} \qquad z = \sin(\theta/2)\frac{\theta_x}{\theta}
\end{equation}
Which means that for every $2\times2$ matrix $U$ in $SU(2)$ one can find $\theta_x, \theta_y, \theta_z$ such that $U = \exp\left[\frac{1}{2}i \theta_i\sigma_i\right]$, and thus $t_i = \frac{1}{2}\sigma_i$ furnish a representation of the $su(2)$ algebra. 
\subsection{The $SO(3)$ group}
The three matrices that make up the representation of the algebra are:
\begin{equation}
    L_1 = \begin{pmatrix}
        0 & 0 & 0 \\
        0 & 0 & -i \\
        0 & i & 0
    \end{pmatrix}
    \qquad
    L_2 = \begin{pmatrix}
        0 & 0 & i \\
        0 & 0 & 0 \\
        -i & 0 & 0
    \end{pmatrix}
    \qquad
    L_3 = \begin{pmatrix}
        0 & -i & 0 \\
        i & 0 & 0 \\
        0 & 0 & 0
    \end{pmatrix}
\end{equation}
Their linear sum with constant $-i\varphi_i$ is
\begin{equation}
    P = (-i) i \begin{pmatrix}
        0 & -\varphi_3 & \varphi_2 \\
        \varphi_3 & 0 & -\varphi_1 \\
        -\varphi_2 & \varphi_1 & 0
    \end{pmatrix} = \begin{pmatrix}
        0 & -\varphi_3 & \varphi_2 \\
        \varphi_3 & 0 & -\varphi_1 \\
        -\varphi_2 & \varphi_1 & 0
    \end{pmatrix}
\end{equation}
Which means that
\begin{equation}
    A = \exp\left[-i\varphi_i L_i\right] = \exp P = \exp \begin{pmatrix}
        0 & -\varphi_3 & \varphi_2 \\
        \varphi_3 & 0 & -\varphi_1 \\
        -\varphi_2 & \varphi_1 & 0
    \end{pmatrix}
\end{equation}
To show that every matrix in $SO(3)$ has the form of $A$, we need to open the exponent. The powers of $P$ are:
\begin{equation}
    P^2 = \begin{pmatrix}
        0 & -\varphi_3 & \varphi_2 \\
        \varphi_3 & 0 & -\varphi_1 \\
        -\varphi_2 & \varphi_1 & 0
    \end{pmatrix}\begin{pmatrix}
        0 & -\varphi_3 & \varphi_2 \\
        \varphi_3 & 0 & -\varphi_1 \\
        -\varphi_2 & \varphi_1 & 0
    \end{pmatrix}
    = \begin{pmatrix}
        -\varphi_2^2 - \varphi_3^2 & \varphi_1\varphi_2 & \varphi_1\varphi_3 \\ \varphi_1\varphi_2 & -\varphi_1^2 - \varphi_3^2 & \varphi_2\varphi_3 \\ \varphi_1\varphi_3 & \varphi_2\varphi_3 & -\varphi_1^2 - \varphi_2^2
    \end{pmatrix}
\end{equation}
\begin{equation}
    P^3 = \begin{pmatrix}
        0 & -\varphi_3 & \varphi_2 \\
        \varphi_3 & 0 & -\varphi_1 \\
        -\varphi_2 & \varphi_1 & 0
    \end{pmatrix} \begin{pmatrix}
        -\varphi_2^2 - \varphi_3^2 & \varphi_1\varphi_2 & \varphi_1\varphi_3 \\ \varphi_1\varphi_2 & -\varphi_1^2 - \varphi_3^2 & \varphi_2\varphi_3 \\ \varphi_1\varphi_3 & \varphi_2\varphi_3 & -\varphi_1^2 - \varphi_2^2
    \end{pmatrix} = -\varphi^2 P
\end{equation}
Then, for $n > 0$
\begin{equation}
    P^{2n} = (i\varphi)^{2n-2} P^{2} \qquad P^{2n+1} = (i\varphi)^{2n}P
\end{equation}
\begin{equation}
\begin{gathered}
    A = \exp P = \sum_{n=0}^{\infty} \frac{P^n}{n!} =
    \sum_{n \ even}\frac{P^n}{n!} + \sum_{n \ odd}\frac{P^n}{n!} = 
    \mathbb{1} + \sum_{n > 0 \ even}\frac{(i\varphi)^{n - 1}}{n!}P^2 + \sum_{n \ odd}\frac{(i\varphi)^{n-1}}{n!}P = 
    \\
    = \mathbb{1} + \sin(\varphi)\frac{P}{\varphi} + \left(1 - \cos\varphi\right)\frac{P^2}{\varphi^2} =
    \\
    = \begin{pmatrix}
        \cos\varphi + n_1^2\left(1 - \cos\varphi\right) &
        n_1n_2\left(1 - \cos\varphi\right) - n_3 \sin\varphi &
        n_1n_3\left(1 - \cos\varphi\right) + n_2 \sin\varphi \\
        n_1n_2\left(1 - \cos\varphi\right) - n_3 \sin(\varphi) &
        \cos(\varphi) + n_2^2\left(1 - \cos(\varphi)\right) &
        n_2n_3\left(1 - \cos\varphi\right) - n_1 \sin(\varphi) \\
        n_1n_3\left(1 - \cos\varphi\right) - n_2 \sin(\varphi) &
        n_2n_3\left(1 - \cos\varphi\right) + n_1 \sin(\varphi) &
        \cos(\varphi) + n_3^2\left(1 - \cos(\varphi)\right)
    \end{pmatrix}
\end{gathered}
\end{equation}
Where $n_i = \varphi_i/\varphi$. This is Rodrigues' rotation formula, which is the most general form of a 3D rotation matrix.
\newline
Now, for any $3\times3$ matrix $O$ we can write:
\begin{equation}
    O = \begin{pmatrix}
        a & b & c \\ d & e & f \\ g & j & h
    \end{pmatrix}
\end{equation} 
And if it's in $SO(3)$, $O^T = O^{-1}$, we can use the generic formula $A^{-1} = \frac{1}{|\det A|}Adj(A)$ and the fact that $\det O = 1$ to get
\begin{equation}
    O^T = O^{-1} = \frac{1}{|\det O|}Adj(O) = Adj(O)
\end{equation}
This gives us the formulas:
\begin{equation}
    \begin{pmatrix} a \\ d \\ g \end{pmatrix} = \begin{pmatrix} eh - fj \\ cj - bh \\ bf - ce \end{pmatrix} = \begin{pmatrix}
        c \\ f \\ h
    \end{pmatrix} \times \begin{pmatrix}
        b \\ e \\ j
    \end{pmatrix}
\end{equation}
\begin{equation}
    \begin{pmatrix} b \\ e \\ j \end{pmatrix} =
    \begin{pmatrix} a \\ d \\ g \end{pmatrix}
    \times
    \begin{pmatrix} c \\ f \\ h \end{pmatrix}
    \qquad
    \begin{pmatrix} b \\ e \\ j \end{pmatrix} =
    \begin{pmatrix} a \\ d \\ g \end{pmatrix}
    \times
    \begin{pmatrix} c \\ f \\ h \end{pmatrix}
\end{equation}
In other words, the three columns of the matrix are an orthonormal basis of $\R^3$ (what a surprise! It was definity not in the group's name or anything). Every orthonormal basis can be displayed as a rotation from some preselected reference base. This means the $SO(3)$ matrices are rotation matrices, and can therefore be written in the form of the Rodrigues formula.

\subsection{Commutation relations}
It's easy to see that for $t_i = \frac{1}{2}\sigma_i$, $\left[t_i,t_j\right] = i\varepsilon_{ijk}t_k$. The simplest way is just to check them one by one:
\begin{equation}
\begin{gathered}
    \left[t_1,t_1\right] = 0
    \qquad
    \left[t_2,t_2\right] = 0
    \qquad
    \left[t_3,t_3\right] = 0
    \\
    \left[t_1,t_2\right] = - \left[t_2,t_1\right] = t_1t_2 - t_2t_1 = \\ = \frac{1}{4}\left[\begin{pmatrix}0&1\\1&0\end{pmatrix}\begin{pmatrix}0&-i\\i&0\end{pmatrix} - \begin{pmatrix}0&-i\\i&0\end{pmatrix}\begin{pmatrix}0&1\\1&0\end{pmatrix}\right] = \frac{1}{4}\left[\begin{pmatrix}i&0\\0&-i\end{pmatrix} - \begin{pmatrix}-i&0\\0&i\end{pmatrix}\right] = \\ = \frac{1}{4}\left[i\sigma_3 - \left(-i\right)\sigma_3\right] = \frac{1}{2}i\sigma_3 = it_3
    \\
    \left[t_1,t_3\right] = - \left[t_3,t_1\right] = t_1t_3 - t_3t_1 = \\ = \frac{1}{4}\left[\begin{pmatrix}0&1\\1&0\end{pmatrix}\begin{pmatrix}1&0\\0&-1\end{pmatrix} - \begin{pmatrix}1&0\\0&-1\end{pmatrix}\begin{pmatrix}0&1\\1&0\end{pmatrix}\right] = \frac{1}{4}\left[\begin{pmatrix}0&-1\\1&0\end{pmatrix} - \begin{pmatrix}0&1\\-1&0\end{pmatrix}\right] = \\ = \frac{1}{4}\left[-i\sigma_2 + \left(-i\right)\sigma_2\right] = -\frac{1}{2}i\sigma_2 = -it_2
    \\
    \left[t_2,t_3\right] = - \left[t_3,t_2\right] = t_2t_3 - t_3t_2 = \\ = \frac{1}{4}\left[\begin{pmatrix}0&-i\\i&0\end{pmatrix}\begin{pmatrix}1&0\\0&-1\end{pmatrix} - \begin{pmatrix}1&0\\0&-1\end{pmatrix}\begin{pmatrix}0&-i\\i&0\end{pmatrix}\right] = \frac{1}{4}\left[\begin{pmatrix}0&i\\i&0\end{pmatrix} - \begin{pmatrix}0&-i\\-i&0\end{pmatrix}\right] = \\ = \frac{1}{4}\left[i\sigma_1 - \left(-i\right)\sigma_1\right] = \frac{1}{2}i\sigma_1 = it_1
\end{gathered}
\end{equation}
Similarly, for $L_i$ we get
\begin{equation}
\begin{gathered}
    \left[L_1,L_1\right] = 0
    \qquad
    \left[L_2,L_2\right] = 0
    \qquad
    \left[L_3,L_3\right] = 0
    \\
    \left[L_1,L_2\right] = - \left[L_2,L_1\right] = L_1L_2 - L_2L_1 = \\ = 
    \begin{pmatrix} 0 & 0 & 0 \\ 0 & 0 & -i \\ 0 & i & 0 \end{pmatrix}
    \begin{pmatrix} 0 & 0 & i \\ 0 & 0 & 0 \\ -i & 0 & 0 \end{pmatrix}
    -
    \begin{pmatrix} 0 & 0 & i \\ 0 & 0 & 0 \\ -i & 0 & 0 \end{pmatrix}
    \begin{pmatrix} 0 & 0 & 0 \\ 0 & 0 & -i \\ 0 & i & 0 \end{pmatrix}
    = \\ =
    \begin{pmatrix} 0 & 0 & 0 \\ -1 & 0 & 0 \\ 0 & 0 & 0 \end{pmatrix}
    -
    \begin{pmatrix} 0 & -1 & 0 \\ 0 & 0 & 0 \\ 0 & 0 & 0 \\ \end{pmatrix}
    = i\begin{pmatrix} 0 & -i & 0 \\ i & 0 & 0 \\ 0 & 0 & 0\end{pmatrix} = iL_3
    \\
    \left[L_1,L_3\right] = - \left[L_3,L_1\right] = L_1L_3 - L_3L_1 = \\ = 
    \begin{pmatrix} 0 & 0 & 0 \\ 0 & 0 & -i \\ 0 & i & 0 \end{pmatrix}
    \begin{pmatrix} 0 & -i & 0 \\ i & 0 & 0 \\ 0 & 0 & 0\end{pmatrix}
    -
    \begin{pmatrix} 0 & -i & 0 \\ i & 0 & 0 \\ 0 & 0 & 0\end{pmatrix}
    \begin{pmatrix} 0 & 0 & 0 \\ 0 & 0 & -i \\ 0 & i & 0 \end{pmatrix}
    = \\ =
    \begin{pmatrix} 0 & 0 & 0 \\ 0 & 0 & 0 \\ -1 & 0 & 0 \end{pmatrix}
    -
    \begin{pmatrix} 0 & 0 & -1 \\ 0 & 0 & 0 \\ 0 & 0 & 0 \\ \end{pmatrix}
    = i\begin{pmatrix} 0 & 0 & -i \\ 0 & 0 & 0 \\ i & 0 & 0\end{pmatrix} = -iL_2
    \\
    \left[L_2,L_3\right] = - \left[L_3,L_2\right] = L_2L_3 - L_3L_2 = \\ = 
    \begin{pmatrix} 0 & 0 & i \\ 0 & 0 & 0 \\ -i & 0 & 0\end{pmatrix}
    \begin{pmatrix} 0 & -i & 0 \\ i & 0 & 0 \\ 0 & 0 & 0\end{pmatrix}
    -
    \begin{pmatrix} 0 & -i & 0 \\ i & 0 & 0 \\ 0 & 0 & 0\end{pmatrix}
    \begin{pmatrix} 0 & 0 & i \\ 0 & 0 & 0 \\ -i & 0 & 0\end{pmatrix}
    = \\ =
    \begin{pmatrix} 0 & 0 & 0 \\ 0 & 0 & 0 \\ 0 & -1 & 0 \end{pmatrix}
    -
    \begin{pmatrix} 0 & 0 & 0 \\ 0 & 0 & -1 \\ 0 & 0 & 0 \\ \end{pmatrix}
    = i\begin{pmatrix} 0 & 0 & 0 \\ 0 & 0 & -i \\ 0 & i & 0\end{pmatrix} = iL_1
\end{gathered}
\end{equation}
This means that both $t_i$ and $L_i$ satisfy the same commutation relations, and there is an obvious isomorphism between their respective algebras.

\subsection{Equivilance of transformations in $SU(2)$ and $SO(3)$}
For every $x_i \in \R^3$, let the $2\times2$ matrix be 
\begin{equation}
    X = x_i\cdot\sigma_i = \begin{pmatrix} x_3 & x_1 - i x_2 \\ x_1 + i x_2 & x_3 \end{pmatrix}
\end{equation}
Now, using equation \ref{eq:su_final_rep}, we can transform $X$ as such:
\begin{equation}
    U = \cos(\theta/2)\mathbb{1} + \sin(\theta/2) \frac{P}{\theta}
    \qquad
    P = \begin{pmatrix}
            i\theta_z & \theta_y + i\theta_x \\
            - \theta_y+ i\theta_x & -i\theta_z
        \end{pmatrix} = i\sum_k \theta_k \sigma_k
\end{equation}
\begin{equation}
\begin{gathered}
    U^\dagger X U = \\ =
    \left[\cos(\theta/2)\mathbb{1} + \sin(\theta/2) \frac{P}{\theta}\right]^\dagger X \left[\cos(\theta/2)\mathbb{1} + \sin(\theta/2) \frac{P}{\theta}\right] = \\ =
    \cos^2(\theta/2)X + \theta^{-1}\sin(\theta/2)\cos(\theta/2)[P^\dagger X - X P] + \theta^{-2}\sin^2(\theta/2)P^\dagger X P
\end{gathered}
\end{equation}
Now we can open each term using $\sigma_i\sigma_j = i\varepsilon_{ijk}\sigma_k + \delta_{ij}\mathbb{1}$ and $\sigma^\dagger = \sigma$:
\begin{equation}
\begin{gathered}
    XP = (x_i\sigma_i)(i\theta_j\sigma_j) = ix_i\theta_j \sigma_i \sigma_j = ix_i\theta_j (i\varepsilon_{ijk}\sigma_k + \delta_{ij}\mathbb{1}) = -x_i\theta_j\varepsilon_{ijk}\sigma_k + ix_i\theta_i\mathbb{1}
    \\
    P^\dagger X = (i\theta_j\sigma_j^\dagger)(x_i\sigma_i) = ix_i\theta_j \sigma_j \sigma_i = ix_i\theta_j (-i\varepsilon_{ijk}\sigma_k + \delta_{ij}\mathbb{1}) = x_i\theta_j\varepsilon_{ijk}\sigma_k + ix_i\theta_i\mathbb{1}
    \\
    P^\dagger X - XP = -2x_i\theta_j\varepsilon_{ijk}\sigma_k = 2(\vec{\theta}\times\vec{x})\cdot\vec{\sigma}
    \\
    P^\dagger X P = (x_i\theta_j\varepsilon_{ijk}\sigma_k + ix_i\theta_i\mathbb{1})(i\theta_l\sigma_l) = ix_i\theta_j\theta_l\varepsilon_{ijk}\sigma_k\sigma_l - x_i\theta_i\theta_l\sigma_l = \\ = 
    ix_i\theta_j\theta_l\varepsilon_{ijk}(i\varepsilon_{klm}\sigma_m + \delta_{kl}\mathbb{1}) - x_i\theta_i\theta_l\sigma_l = \\ =
    -x_i\theta_j\theta_l\varepsilon_{ijk}\varepsilon_{klm}\sigma_m + ix_i\theta_j\theta_l\varepsilon_{ijk}\delta_{kl}\mathbb{1} - x_i\theta_i\theta_l\sigma_l = \\ =
    -x_i\theta_j\theta_l(\delta_{il}\delta_{jm} - \delta_{im}\delta_{jl})\sigma_m + ix_i\theta_j\theta_l\varepsilon_{ijl}\mathbb{1} - x_i\theta_i\theta_l\sigma_l
\end{gathered}
\end{equation}
The third term this equaiton is a product of a symmetric tensor $x_i\theta_j\theta_l$ and an antisymmetric tensor $\varepsilon_{ijl}$, which means it's zero. This leaves:
\begin{equation}
\begin{gathered}
    P^\dagger X P = -x_i\theta_j\theta_i\sigma_j + x_i\theta_j\theta_j\sigma_i - x_i\theta_i\theta_l\sigma_l = \theta^2\vec{x}\cdot\vec{\sigma} - 2(\vec{x}\cdot\vec{\theta})(\vec{\theta}\cdot\vec{\sigma})
\end{gathered}
\end{equation}
In total, we get
\begin{equation}
\begin{gathered}
    U^\dagger X U = \\ = \cos^2(\theta/2)X + \theta^{-1}\sin(\theta/2)\cos(\theta/2)\cdot2(\vec{\theta}\times\vec{x})\cdot\vec{\sigma} + \theta^{-2}\sin^2(\theta/2)\left(\theta^2\vec{x}\cdot\vec{\sigma} - 2(\vec{x}\cdot\vec{\theta})(\vec{\theta}\cdot\vec{\sigma})\right) = \\ =
    \cos^2(\theta/2)X + 2\theta^{-1}\sin(\theta/2)\cos(\theta/2)(\vec{\theta}\times\vec{x})\cdot\vec{\sigma} + \sin^2(\theta/2)X - 2\theta^{-2}\sin^2(\theta/2)(\vec{x}\cdot\vec{\theta})(\vec{\theta}\cdot\vec{\sigma}) = \\ =
    X + \theta^{-1}\sin(\theta)(\vec{\theta}\times\vec{x})\cdot\vec{\sigma} - \theta^{-2}(1 - \cos(\theta))(\vec{\theta}\cdot\vec{x})(\vec{\theta}\cdot\vec{\sigma}) = \\ =
    \left(x_i + \theta^{-1}\sin(\theta)(\varepsilon_{jki} \theta_j x_k) - \theta^{-2}(1 - \cos(\theta))\theta_j\theta_ix_j\right)\sigma_i
\end{gathered}
\end{equation}
defining $L_{ij} = \varepsilon_{kij}\theta_k$, we get:
\begin{equation}
    L^2x = L\times(L\times x) = L(Lx) = \theta_i\theta_jx_j
\end{equation}
Which means
\begin{equation}
    U^\dagger X U = \left(x_i + \theta^{-1}\sin(\theta)(\varepsilon_{jki} \theta_j x_k) - \theta^{-2}(1 - \cos(\theta))\theta_j\theta_ix_j\right)\sigma_i
\end{equation}
Which is exactly the Rodrigues Rotation formula, meanign the vector $x_i$ transforms like a vector in 3d under an $SO(3)$ transformation.
\newline
This is the $SU(2) \rightarrow SO(3)$ mapping
\begin{equation}
    U \in SU(2) \sim O \in SO(3) \ if \ U^\dagger X U = (O\vec{x})\cdot\vec{\sigma}
\end{equation}
Since $U$ appears twice in the equation, the same matrix $O$ maps to both $U$ and $-U$, making this mapping a surjective two-to-one homomorphism.
\newpage
\section{A bit of $SO(1,3)$}
\subsection{Weyl Spinor and Lorentz transformation relations}
Using the hist, let us first show that any Pauli matrix $\sigma_n$ satisfies $\sigma_n\sigma_2 = -\sigma_2\sigma_n^T$. The easiest way to show this is just to calculte both sides for each of the matrices and see that they give the same result. We will use the identities $\sigma_i\sigma_j = i\varepsilon_{ijk}\sigma_k + \delta_{ij}\mathbb{1}$ and $\sigma_1^T = \sigma_1, \sigma_2^T = -\sigma_2,$ and $\sigma_3^T = \sigma_3$.
\begin{equation}
\begin{gathered}
    \sigma_1\sigma_2 = i\sigma_3 \qquad -\sigma_2\sigma_1^T = -\sigma_2\sigma_1 = -(-i\sigma_3) = i\sigma_3
    \\
    \sigma_2\sigma_2 = \mathbb{1} \qquad -\sigma_2\sigma_2^T = \sigma_2\sigma_2 = \mathbb{1}
    \\
    \sigma_3\sigma_2 = -i\sigma_1 \qquad -\sigma_2\sigma_3^T = -\sigma_2\sigma_3 = -(i\sigma_1) = -i\sigma_1
\end{gathered}
\end{equation}
This property is true, not only for $\sigma_n$, but also for $\sigma_n^T$, since they are just $\sigma_n$ up to a constant.
\newline
Using the representation we found in \ref{eq:su_final_rep}:
\begin{equation}
\begin{gathered}
    M_L^{T}\sigma_2 = \exp(-i\left(\theta_n - iy_n\right)\frac{\sigma_n}{2})^T\sigma_2 =
    \left[\cos(\frac{\sqrt{\theta^2 + y^2}}{2})\mathbb{1} - \sin(\frac{\sqrt{\theta^2 + y^2}}{2}) \frac{\theta_n - iy_n}{\sqrt{\theta^2 + y^2}}\sigma_n^T\right]\sigma_2
    = \\ =
    \left[\cos(\frac{\sqrt{\theta^2 + y^2}}{2})\sigma_2 + \sin(\frac{\sqrt{\theta^2 + y^2}}{2}) \frac{\theta_n - iy_n}{\sqrt{\theta^2 + y^2}}\sigma_2\sigma_n\right]
    = \\ =
    \sigma_2\left[\cos(\frac{\sqrt{\theta^2 + y^2}}{2})\mathbb{1} + \sin(\frac{\sqrt{\theta^2 + y^2}}{2}) \frac{\theta_n - iy_n}{\sqrt{\theta^2 + y^2}}\sigma_n\right]
    = \\ =
    \sigma_2\exp(i\left(\theta_n - iy_n\right)\frac{\sigma_n}{2})
    = \sigma_2M_R^\dagger
\end{gathered}
\end{equation}
It is also easy to see that
\begin{equation}
    M_R^\dagger M_L = \exp(i\left(\theta_n - iy_n\right)\frac{\sigma_n}{2})\exp(-i\left(\theta_n - iy_n\right)\frac{\sigma_n}{2}) = \mathbb{1}
\end{equation}
Which is to say $M_R^\dagger = M_L^{-1}$ and so
\begin{equation}
    M_L^T\sigma_2 = \sigma_2M_R^\dagger = \sigma_2M_L^{-1}
\end{equation}
\subsection{Lorentz scalar from the spinor fields}
The Lorentz transformaiton for the vector $A^\mu$ is $A^\mu \rightarrow A'^\mu = \Lambda^\mu_\nu A^\nu$. We saw in section 1 that this is equivilant to transform the matrix $A^\mu \bar{\sigma}_\mu \rightarrow M_R A^\mu \bar{\sigma}_\mu M_R^\dagger$. We saw in class that $\bar{\sigma}_\mu$ transforms using $U = M_R^\dagger$. Using this, together with the transformations given for the spinor fields in the question, we get
\begin{equation}
    \xi^\dagger A^\mu \bar{\sigma}_\mu \xi \rightarrow
    \xi^\dagger M_L^\dagger M_R A^\mu \bar{\sigma}_\mu M_R^\dagger M_L\xi = \xi^\dagger A^\mu \bar{\sigma}_\mu \xi
\end{equation}
There is, of course, a simialr scalar for the $\chi$ field, simply by replacing $\xi \rightarrow \chi, \bar{\sigma}_\mu \rightarrow \sigma_\mu$. Then, $\sigma_\mu \rightarrow M_L \sigma_mu M_L^\dagger$:
\begin{equation}
    \chi^\dagger A^\mu \sigma_\mu \chi \rightarrow
    \chi^\dagger M_R^\dagger M_L A^\mu \sigma_\mu M_L^\dagger M_R\chi = \chi^\dagger A^\mu \sigma_\mu \chi
\end{equation}
In both cases we used $M_L^{-1} = M_R^\dagger$.

\section{What adds under Lorentz transformations?}
\subsection{Rotations}
For a rotation around the z-axis, the $4\times4$ Lorentz transformation is
\begin{equation}
    \Lambda^\mu_\nu(\theta) = \begin{pmatrix}
        1 & 0 & 0 & 0 \\
        0 & \cos\theta & -\sin\theta & 0 \\
        0 & \sin\theta & \cos\theta & 0 \\
        0 & 0 & 0 & 1
    \end{pmatrix}
\end{equation}
Using the well known trigonometric identities for angle addition, it's easy to see that performing two successive rotations around the z-axis by two angles $\theta_1$ and $\theta_2$ is equivilant to a single rotation by $\theta_1 + \theta_2$
\begin{equation}
\begin{gathered}
    \Lambda^\mu_\alpha(\theta_1) \Lambda^\alpha_\nu(\theta_2) =
    \\
    = \begin{pmatrix}
        1 & 0 & 0 & 0 \\
        0 & \cos\theta_1 & -\sin\theta_1 & 0 \\
        0 & \sin\theta_1 & \cos\theta_1 & 0 \\
        0 & 0 & 0 & 1
    \end{pmatrix} \begin{pmatrix}
        1 & 0 & 0 & 0 \\
        0 & \cos\theta_2 & -\sin\theta_2 & 0 \\
        0 & \sin\theta_2 & \cos\theta_2 & 0 \\
        0 & 0 & 0 & 1
    \end{pmatrix} = 
    \\
    = \begin{pmatrix}
        1 & 0 & 0 & 0 \\
        0 & \cos\theta_1\cos\theta_2 - \sin\theta_1\sin\theta_2 & -\cos\theta_1\sin\theta_2-\sin\theta_1\cos\theta_2 & 0 \\
        0 & \sin\theta_1\cos\theta_2 + \cos\theta_1\sin\theta_2 & \sin\theta_1\sin\theta_2 + \cos\theta_1\cos\theta_2 & 0 \\
        0 & 0 & 0 & 1
    \end{pmatrix} = 
    \\
    = \begin{pmatrix}
        1 & 0 & 0 & 0 \\
        0 & \cos\left(\theta_1+\theta_2\right) & -\sin\left(\theta_1+\theta_2\right) & 0 \\
        0 & \sin\left(\theta_1+\theta_2\right) & \cos\left(\theta_1+\theta_2\right) & 0 \\
        0 & 0 & 0 & 1
    \end{pmatrix} = \Lambda^\mu_\nu(\theta_1 + \theta_2)
\end{gathered}
\end{equation}
\subsection{Boosts}
A boost along the x-axis by a velocity parameter $\beta$ is
\begin{equation}
    \Lambda^\mu_\nu(\beta) = \begin{pmatrix}
        \gamma & -\beta\gamma & 0 & 0 \\
        -\beta\gamma & \gamma & 0 & 0 \\\
        0 & 0 & 1 & 0 \\
        0 & 0 & 0 & 1
    \end{pmatrix}
\end{equation}
Where $\gamma = \left(1-\beta^2\right)^{-1/2}$. Performing two consecutive such transformations by two different velocities $\beta_1$ and $\beta_2$, will amount to
\begin{equation}
\begin{gathered}
    \Lambda^\mu_\alpha(\beta_1) \Lambda^\alpha_\nu(\beta_2) =
    \\
    = \begin{pmatrix}
        \gamma_1 & -\beta_1\gamma_1 & 0 & 0 \\
        -\beta_1\gamma_1 & \gamma_1 & 0 & 0 \\
        0 & 0 & 1 & 0 \\
        0 & 0 & 0 & 1
    \end{pmatrix}
    \begin{pmatrix}
        \gamma_2 & -\beta_2\gamma_2 & 0 & 0 \\
        -\beta_2\gamma_2 & \gamma_2 & 0 & 0 \\
        0 & 0 & 1 & 0 \\
        0 & 0 & 0 & 1
    \end{pmatrix} =
    \\
    = \begin{pmatrix}
        \gamma_1\gamma_2\left(1 + \beta_1\beta_2\right) & -\gamma_1\gamma_2\left(\beta_1 + \beta_2\right) & 0 & 0 \\
        -\gamma_1\gamma_2\left(\beta_1 + \beta_2\right) & \gamma_1\gamma_2\left(1 + \beta_1\beta_2\right) & 0 & 0 \\
        0 & 0 & 1 & 0 \\
        0 & 0 & 0 & 1
    \end{pmatrix}
\end{gathered}
\end{equation}
Let us now look for what $\beta$ and $\gamma$ fit this new transformation. This means we are looking for $\beta$ and its matching $\gamma$ for which $\Lambda^\mu_\alpha(\beta_1)\Lambda^\alpha_\nu(\beta_2) = \Lambda^\mu_\nu(\beta)$, so we can find the addative connection. From the last matrix, we can see that
\begin{equation}
    \gamma = \gamma_1\gamma_2(1 + \beta_1\beta_2) \qquad \-\beta\gamma = -\gamma_1\gamma_2(\beta_1 + \beta_2)
\end{equation}
Therefore
\begin{equation}
\label{eq:betas_relation}
    \beta = \frac{\gamma\beta}{\gamma} = \frac{\gamma_1\gamma_2(\beta_1 + \beta_2)}{\gamma_1\gamma_2(1+\beta_1\beta_2)} = \frac{\beta_1+\beta_2}{1+\beta_1\beta_2}
\end{equation}
This reminds us of the addition formula for hyperbolic tangent:
\begin{equation}
    \tanh(\alpha+\beta) = \frac{\tanh\alpha + \tanh\beta}{1 + \tanh\alpha\tanh\beta}
\end{equation}
So let us define the rapidity, with a range $(-\infty,\infty)$:
\begin{equation}
    y(\beta) = \tanh^{-1}(\beta)
\end{equation}
and thus equation \ref{eq:betas_relation} becomes
\begin{equation}
    \tanh(y) = \frac{\tanh(y_1) + \tanh(y_2)}{1 + \tanh(y_1)\tanh(y_2)}
\end{equation}
\begin{equation}
    y = \tanh^{-1}\left(\frac{\tanh(y_1) + \tanh(y_2)}{1 + \tanh(y_1)\tanh(y_2)}\right) = \tanh^{-1}(\tanh(y_1 + y_2)) = y_1 + y_2
\end{equation}
Which means
\begin{equation}
    \Lambda^\mu_\alpha(y_1)\Lambda^\alpha_\nu(y_2) = \Lambda^\mu_\nu(y_1 + y_1)
\end{equation}

\section{Translation in field space}
\subsection{Translation using the Noether charge}
Opening the equation using Campbell identity
\begin{equation}
    e^{i\alpha Q}\phi(x)e^{-i\alpha Q} = \sum_{n=0}^\infty \frac{1}{n!}\left[i\alpha Q, \phi(x)\right]_n
\end{equation}
Where
\begin{equation}
    \left[A, B\right]_n = [A, [A,B]_{n-1}] \qquad [A,B]_0 = B
\end{equation}
Specifically, for the charge $Q = \int d^3x \pi$, we can use the commutation relations
\begin{equation}
    \left[\phi(x),\pi(y)\right] = i\delta^3(x - y)
\end{equation}
Which means
\begin{equation}
    [i\alpha Q, \phi(x)] = i\alpha \int d^3y [\pi(y), \phi(x)] = i\alpha \int d^3y (-i)\delta^3(x - y) = \alpha
\end{equation}
And so every commutation relations after the second order vanish, and we are left with
\begin{equation}
    e^{i\alpha Q}\phi(x)e^{-i\alpha Q} = \phi(x) + \alpha
\end{equation}
\subsection{Additional Lorentz invariants}
Any other Lorentz invariant that is only dependent on the derivative of the field is also invariant to an addition of a constant to the field. The most simple function of derivatives is polynomials, and the next polynomial that is not already in the Lagrangian is the quadruple derivative:
\begin{equation}
    S = \int d^4x \frac{1}{2}\partial_\mu\phi\partial^\mu\phi + \frac{1}{4}\left(\partial_\mu\phi\partial^\mu\phi\right)^2
\end{equation}

\subsection{The new Conjugate momentum}
The new conjugate momentum is
\begin{equation}
    \pi = \frac{\partial \mathcal{L}}{\partial \dot{\phi}} = [1 + \left(\partial_\mu\phi\partial^\mu\phi\right)]\dot{\varphi}
\end{equation}
Using this new momentum, we can see if the shift in the field space is stil the same. For this we will calculate the commutation relations:
\begin{equation}
    [\phi(x),\pi(y)] = [\dot{\varphi} + \partial^{(x)}_\mu\phi\partial^{(x)\mu}\phi\dot{\varphi}, pi(y)] = i\delta^3(x-y)(1 + \partial^{(x)}_\mu \partial^{(x)\mu})
\end{equation}
And thus, since the derivative of 1 vanishes, we get:
\begin{equation}
    [i\alpha Q, \phi(x)] = i\alpha \int d^3y [\pi(y), \phi(x)] = i\alpha \int d^3y (-i)\delta^3(x - y)(1 + \partial^{(x)}_\mu \partial^{(x)\mu}) = \alpha
\end{equation}
Meaning that once again, the Campbell identity gets cut off after the linear term in $\alpha$, and the momentum once again shifts the field by a constant.

\section{The Baker-Campbell-Hausdorff formula}
\subsection{Proving for the first and second order}
Opening the exponents up to second order:
\begin{equation}
\begin{gathered}
    e^Ae^B = \left(1 + A + \frac{1}{2}A^2 + O(A^3)\right)\left(1 + B + \frac{1}{2}B^2 + O(B^3)\right) = \\ =
    1 + A + B + \frac{1}{2}A^2 + AB + \frac{1}{2}B^2 + \cdots = \\ =
    1 + (A + B) + \frac{1}{2}(A + B)^2 + \frac{1}{2}[A,B] + \cdots = \\ = \exp(A + B + \frac{1}{2}[A,B] + \cdots) = e^C
\end{gathered}
\end{equation}
As for the third order, there each term requires a different commutator, since the third term in the exponent of $A+B$ is
\begin{equation}
    \frac{1}{6}(A+B)^3 = \frac{1}{6}\left(A^3 + A^2B + AB^2 + B^3 + B^2A + BA^2 + ABA + BAB\right)
\end{equation}
While the third order of the result from the product of the exponent $e^Ae^B$ is
\begin{equation}
    \frac{1}{2}AB^2 + \frac{1}{2}A^2B + \frac{1}{6}A^3 + \frac{1}{6}B^3
\end{equation}
Subtracting them will give us the remainder, which is the third order term in $C$:
\begin{equation}
\begin{gathered}
    \frac{1}{2}AB^2 + \frac{1}{2}A^2B + \frac{1}{6}A^3 + \frac{1}{6}B^3 - \frac{1}{6}\left(A^3 + A^2B + AB^2 + B^3 + B^2A + BA^2 + ABA + BAB\right) = \\ =
    \frac{1}{6}\left(2AB^2 - A^2B - ABA + 2BA^2 - B^2A - BAB\right) = \\ =
    \frac{1}{6}\left([AB, B] + [BA, A] + [A,B^2] + [B,A^2]\right)
\end{gathered}
\end{equation}
Which means that the third order of the approximation is also made of commutators.
\subsection{The $SL(2,C)$ group and the divergence of the identity}
The gropu $SL(2,C)$ is defined as the group of all complex $2x2$ matrices with determinant 1. Now, let us take a matrix $S \in SL(2, C)$ very close to the identity:
\begin{equation}
    S \approx \mathbb{1} + \epsilon = \begin{pmatrix}
        1 + \epsilon_{11} & \epsilon_{12} \\ \epsilon_{21} & 1 + \epsilon_{22}
    \end{pmatrix}
\end{equation}
This matrix must satisfy $\det S = 1$, and so:
\begin{equation}
    (1 + \epsilon_{11})(1 + \epsilon_{22}) - \epsilon_{12}\epsilon_{21} = 1 + \epsilon_{11} + \epsilon_{22} + \epsilon_{11}\epsilon_{22} - \epsilon_{12}\epsilon_{21} = 1 + \tr(\epsilon) + \det(\epsilon) = 1
\end{equation}
Since $\det \epsilon$ is of the second order, we can neglect it and get
\begin{equation}
    \tr \epsilon = 0
\end{equation}
Or, in other words
\begin{equation}
    \epsilon_{11} = -\epsilon_{22}
\end{equation}
Which means that a simple basis of three traceless matrices for the algebra is
\begin{equation}
    A = \begin{pmatrix}1 & 0 \\ 0 & -1\end{pmatrix}
    \qquad
    B = \begin{pmatrix}0 & 1 \\ 0 & 0\end{pmatrix}
    \qquad
    C = \begin{pmatrix}0 & 0 \\ 1 & 0\end{pmatrix}
\end{equation}
Now we can look at the matrices $X$ and $Y$ from $SL(2, C)$:
\begin{equation}
    X = \begin{pmatrix}0 & i\pi \\ i\pi & 0\end{pmatrix}
    \qquad
    Y = \begin{pmatrix}0 & -1 \\ 0 & 0\end{pmatrix}
\end{equation}
To calcualte their exponents, we will first calcualte their powers:
\begin{equation}
    X^2 = -\pi^2 \mathbb{1} \qquad Y^2 = 0
\end{equation}
Therefore:
\begin{equation}
\begin{gathered}
    e^X = \sum_n \frac{X^n}{n!} = \sum_{n \ even} \frac{X^n}{n!} + \sum_{n \ odd} \frac{X^n}{n!} = \sum_{n \ even} \frac{(i\pi)^n}{n!}\mathbb{1} + \sum_{n \ odd} \frac{(i\pi)^{n-1}}{n!}X = \\ =
    \cos(\pi)\mathbb{1} + \frac{1}{\pi}\sin(\pi)X = -\mathbb{1}
\end{gathered}
\end{equation}
Which is, of course, part of the group, and
\begin{equation}
    e^{Y} = \sum_n \frac{X^n}{n!} = \mathbb{1} + Y = \begin{pmatrix}
        1 & -1 \\ 0 & 1
    \end{pmatrix}
\end{equation}
which is also a part of the group. However:
\begin{equation}
    e^Xe^Y = \begin{pmatrix}
        -1 & 0 \\ 0 & -1
    \end{pmatrix}\begin{pmatrix}
        1 & -1 \\ 0 & 1
    \end{pmatrix} = \begin{pmatrix}
        -1 & 1 \\ 0 & -1
    \end{pmatrix}
\end{equation}
The determinant of which is $\det(e^Xe^Y) = 1$, which means it is a part of $SL(2,C)$. However, if we try to find the matrix $Z$ such that $e^Z = e^Xe^Y$, we will find that we can't. The proof goes as follows:
\newline
The matrix $Z$ has two eigenvalues which could be differen or equal. If they are different, they must satisfy $\lambda_1 = -\lambda_2$, since $\tr Z = 0$. This would mean $Z$ is diagonalizable and thus $e^Z$ is also diagonalizable. If the two eigenvalues are equal, then they must be zero, which means that there exist a vector $\vec{v}$ such that $e^Z\vec{v} = e^0\vec{v} = 1\vec{v} = \vec{v}$, meaning that $1$ must be an eigenvalue of the matrix $e^Z$. We got a conclusion that every matrix that can be displayed as $e^Z$ is either diagonalizable or has an eigenvalue $1$. Since the matrix that we got is doesn't satisfy any of these requirements, it canno't be represented in this form. 
\end{document}